\Needspace{8\baselineskip}
\Article{177}
\textsc{Aufgabe.} \emph{Gegeben ist eine Folge von Formen, in der jede ihrer Vorgängerin nach rechts benachbart ist. Gesucht ist eine eigentliche Transformation der ersten Form in eine beliebige Form der Folge.}

\SourceLeaf{158}
\emph{Lösung.} Die Formen seien $(a,b,a')=F$; $(a',b',a'')=F'$; $(a'',b'',a''')=F''$; $(a''',b''',a'''')=F'''$ usw. Wir bezeichnen $\frac{b+b'}{a'}$, $\frac{b'+b''}{a''}$, $\frac{b''+b'''}{a'''}$ usw. jeweils mit $h',h'',h'''$ usw. Die Unbestimmten der Formen $F,F',F''$ usw. seien $x,y$; $x',y'$; $x'',y''$ usw. Wir nehmen an, dass $F$ übergeht
\[
\begin{array}{ll@{\qquad}l@{\qquad}l}
\text{in }F'&\text{bei}&x=\alpha'x'+\beta'y',&y=\gamma'x'+\delta'y'\\
F''&\ldots&x=\alpha''x''+\beta''y'',&y=\gamma''x''+\delta''y''\\
F'''&\ldots&x=\alpha'''x'''+\beta'''y''',&y=\gamma'''x'''+\delta'''y'''\\
&&\text{usw.}&
\end{array}
\]
{\small\setlength{\arraycolsep}{3pt}
\[
\begin{array}{l@{\enspace}l@{\enspace}l}
\text{Da nun }F\text{ in }F'\text{ übergeht bei}&x=-y',&y=x'+h'y'\\
\hfill F'\text{ in }F''\text{ bei}&x'=-y'',&y'=x''+h''y''\\
\hfill F''\text{ in }F'''\text{ bei}&x''=-y''',&y''=x'''+h'''y'''\text{ usw. (Artikel\ 160)}
\end{array}
\]
}
so ergibt sich leicht folgender Algorithmus (Artikel 159):
\[
\begin{array}{ll@{\qquad}ll}
\alpha'=0&\beta'=-1&\gamma'=1&\delta'=h'\\
\alpha''=\beta'&\beta''=h''\beta'-\alpha'&\gamma''=\delta'&\delta''=h''\delta'-\gamma'\\
\alpha'''=\beta''&\beta'''=h'''\beta''-\alpha''&\gamma'''=\delta''&\delta'''=h'''\delta''-\gamma''\\
\alpha''''=\beta'''&\beta''''=h''''\beta'''-\alpha'''&\gamma''''=\delta'''&\delta''''=h''''\delta'''-\gamma'''\\
&\text{usw.}&&
\end{array}
\]
oder
\[
\begin{array}{ll@{\qquad}ll}
\alpha'=0&\beta'=-1&\gamma'=1&\delta'=h'\\
\alpha''=\beta'&\beta''=h''\beta'&\gamma''=\delta'&\delta''=h''\delta'-1\\
\alpha'''=\beta''&\beta'''=h'''\beta''-\beta'&\gamma'''=\delta''&\delta'''=h'''\delta''-\delta'\\
\alpha''''=\beta'''&\beta''''=h''''\beta'''-\beta''&\gamma''''=\delta'''&\delta''''=h''''\delta'''-\delta''\\
&\text{usw.}&&
\end{array}
\]

Dass alle diese Transformationen eigentlich sind, folgt ohne Mühe sowohl aus ihrer Konstruktion als auch aus Artikel 159.

Dieser sehr einfache und für die Rechnung praktische Algorithmus ist dem in Artikel 27 beschriebenen Algorithmus analog und lässt sich auch auf ihn zurückführen\HistoricalNote{*)}{In der Bezeichnung von Artikel 27 gilt nämlich $\beta^n=\pm[-h'',h''',-h''''\ldots\pm h^n]$, wobei für die jeweils zur Wahl stehenden Vorzeichen $-\,-$; $-\,+$; $+\,-$; $+\,+$ zu nehmen sind, je nachdem, ob $n$ die Form $4k+0$; $1$; $2$; $3$ hat, und
\[
\delta^n=\pm[h',-h'',h'''\ldotp\ldotp\ldotp\ldotp\ldotp\pm h^n]
\]
wobei die zur Wahl stehenden Vorzeichen $+\,-$; $+\,+$; $-\,-$; $-\,+$ sein müssen, je nachdem, ob $n$ die Form $4k+0$; $1$; $2$; $3$ hat. Um uns kurz zu halten, können wir dies hier nicht ausführlicher entwickeln; jeder kann es aber leicht selbst nachprüfen.}. Diese Lösung
\SourceLeaf{159}
ist nicht auf Formen mit negativer Determinante beschränkt, sondern gilt in allen Fällen, sofern keine der Zahlen $a',a'',a'''$ usw. $=0$ ist.

\Needspace{8\baselineskip}
\Article{178}
\textsc{Aufgabe.} \emph{Zu zwei gegebenen, eigentlich äquivalenten Formen $F$, $f$ mit derselben negativen Determinante eine eigentliche Transformation der einen in die andere finden.}

\emph{Lösung.} Wir nehmen an, dass die Form $F$ gleich $(A,B,A')$ ist und man mit der Methode von Artikel 171 die Folge der Formen $(A',B',A'')$, $(A'',B'',A''')$ usw. bis zu $(A^m,B^m,A^{m+1})$ gefunden hat, wobei die letzte Form reduziert ist. Ebenso sei $f$ gleich $(a,b,a')$, und mit derselben Methode sei die Folge $(a',b',a'')$, $(a'',b'',a''')$ bis zu $(a^n,b^n,a^{n+1})$ gefunden worden, wobei die letzte Form reduziert ist. Dann sind zwei Fälle möglich.

\begingroup\rightskip=0pt plus 3em\relax
I. Die Formen $(A^m,B^m,A^{m+1})$, $(a^n,b^n,a^{n+1})$ sind entweder identisch oder entgegengesetzt und zugleich ambig. Dann sind die Formen $(A^{m-1},B^{m-1},A^m)$, $(a^n,-b^{n-1},a^{n-1})$ benachbart (dabei bezeichnet $A^{m-1}$ den vorletzten Term der Folge $A,A',A''\ldots A^m$, und $B^{m-1}$, $a^{n-1}$, $b^{n-1}$ haben entsprechende Bedeutungen). Denn $A^m=a^n$, $B^{m-1}\equiv-B^m\ (\mathrm{mod.}\ A^m)$, $b^{n-1}\equiv-b^n\ (\mathrm{mod.}\ a^n\text{ oder }A^m)$, also $B^{m-1}-b^{n-1}\equiv b^n-B^m$ und damit $\equiv0$, wenn die Formen $(A^m,B^m,A^{m+1})$, $(a^n,b^n,a^{n+1})$ identisch sind, und $\equiv2b^n$, also $\equiv0$, wenn sie entgegengesetzt und ambig sind. Somit ist in der Folge der Formen
\[
\begin{gathered}
(A,B,A'),\ (A',B',A'')\ldotp\ldotp\ldotp\ldotp(A^{m-1},B^{m-1},A^m),\\
(a^n,-b^{n-1},a^{n-1}),\ (a^{n-1},-b^{n-2},a^{n-2})\ldots(a',-b,a),\ (a,b,a')
\end{gathered}
\]
jede Form ihrer Vorgängerin benachbart. Nach dem vorigen Artikel lässt sich also eine eigentliche Transformation der ersten Form $F$ in die letzte $f$ finden.
\par\endgroup

II. Die Formen $(A^m,B^m,A^{m+1})$, $(a^n,b^n,a^{n+1})$ sind nicht identisch, sondern entgegengesetzt, und zugleich gilt $A^m=A^{m+1}=a^n=a^{n+1}$. Dann besitzt die Folge der Formen
\[
\begin{gathered}
(A,B,A'),\ (A',B',A'')\ldotp\ldotp\ldotp\ldotp(A^m,B^m,A^{m+1}),\\
(a^n,-b^{n-1},a^{n-1}),\ (a^{n-1},-b^{n-2},a^{n-2})\ldotp\ldotp\ldotp\ldotp(a',-b,a),\ (a,b,a')
\end{gathered}
\]
dieselbe Eigenschaft. Denn $A^{m+1}=a^n$, und $B^m-b^{n-1}=-(b^n+b^{n-1})$ ist durch $a^n$ teilbar. Nach dem vorigen Artikel findet man also eine eigentliche Transformation der ersten Form $F$ in die letzte $f$.

\SourceLeaf{160}
\emph{Beispiel.} So ergibt sich für die Formen $(23.\,38,63)$, $(15,20,27)$ die Folge
\[
\begin{gathered}
(23,38,63).\ (63,25,10),\ (10,5,3),\ (3,1.\,2),\\
(2.\,-7,27),\ (27,-20,15),\ (15,20,27)
\end{gathered}
\]
und damit
\[
h'=1,\quad h''=3,\quad h'''=2,\quad h''''=-3,\quad h'''''=-1,\quad h''''''=0.
\]
Daraus folgt diese Transformation der Form $23xx+76xy+63yy$ in $15tt+40tu+27uu$: $x=-13t-18u$, $y=8t+11u$.

\Needspace{12\baselineskip}
Aus dieser Lösung folgt ohne Mühe die Lösung der Aufgabe: \emph{Sind die Formen $F,f$ uneigentlich äquivalent, so soll eine uneigentliche Transformation der Form $F$ in $f$ gefunden werden.} Sei nämlich $f=att+2btu+a'uu$. Dann ist die entgegengesetzte Form $app-2bpq+a'qq$ der Form $F$ eigentlich äquivalent. Man sucht eine eigentliche Transformation der Form $F$ in jene, $x=\alpha p+\beta q$, $y=\gamma p+\delta q$. Offenbar geht dann $F$ in $f$ über, wenn man $x=\alpha t-\beta u$, $y=\gamma t-\delta u$ setzt, und diese Transformation ist uneigentlich.

Sind die Formen $F,f$ also sowohl eigentlich als auch uneigentlich äquivalent, so lässt sich sowohl eine eigentliche als auch eine uneigentliche Transformation finden.

\Needspace{8\baselineskip}
\Article{179}
\textsc{Aufgabe.} \emph{Sind die Formen $F$, $f$ äquivalent, so sollen alle Transformationen der Form $F$ in $f$ gefunden werden.}

\emph{Lösung.} Sind die Formen $F$, $f$ nur auf eine Art äquivalent, \emph{das heißt} nur eigentlich oder nur uneigentlich, so sucht man nach dem vorigen Artikel eine Transformation der Form $F$ in $f$. Offenbar kann es nur weitere Transformationen geben, die dieser gleichartig sind. Sind dagegen die Formen $F$, $f$ sowohl eigentlich als auch uneigentlich äquivalent, so sucht man zwei Transformationen, eine eigentliche und eine uneigentliche. Sei nun die Form $F=(A,B,C)$, ferner $BB-AC=-D$ und der größte gemeinsame Teiler der Zahlen $A,2B,C$ $=m$. Aus Artikel 162 folgt: Im ersten Fall lassen sich alle Transformationen der Form $F$ in $f$ aus einer Transformation herleiten; im zweiten Fall alle eigentlichen aus der eigentlichen und alle uneigentlichen aus der uneigentlichen, sofern man alle Lösungen der Gleichung $tt+Duu=mm$ kennt. Sind diese gefunden, ist die Aufgabe also gelöst.

Es gilt aber $D=AC-BB$, $4D=4AC-4BB$, also ist $\frac{4D}{mm}=4\frac{AC}{mm}-(\frac{2B}{m})^2$ eine ganze Zahl. Ist nun

1) $\frac{4D}{mm}>4$, so gilt $D>mm$. Daher muss in $tt+Duu=mm$ notwendig $u$ $=0$ sein, und $t$ kann nur die Werte $+m$ und $-m$ haben. Somit gilt:
\SourceLeaf{161}
Sind $F$, $f$ nur auf eine Art äquivalent, und ist eine Transformation
\[
x=\alpha x'+\beta y',\quad y=\gamma x'+\delta y'
\]
gegeben, so gibt es außer dieser, die aus $t=m$ entsteht (Artikel 162), und der folgenden
\[
x=-\alpha x'-\beta y',\quad y=-\gamma x'-\delta y'
\]
keine weiteren. Sind dagegen $F$, $f$ sowohl eigentlich als auch uneigentlich äquivalent, und hat man eine eigentliche Transformation
\[
x=\alpha x'+\beta y',\quad y=\gamma x'+\delta y'
\]
und eine uneigentliche
\[
x=\alpha'x'+\beta'y',\quad y=\gamma'x'+\delta'y'
\]
so gibt es außer der ersten (aus $t=m$) und dieser
\[
x=-\alpha x'-\beta y',\ y=-\gamma x'-\delta y'\ (\text{aus }t=-m)
\]
keine weitere eigentliche; ebenso gibt es keine weitere uneigentliche außer
\[
\begin{gathered}
x=\alpha'x'+\beta'y',\quad y=\gamma'x'+\delta'y';\\
\text{und }x=-\alpha'x'-\beta'y',\quad y'=-\gamma'x'-\delta'y'
\end{gathered}
\]

2) Ist $\frac{4D}{mm}=4$ oder $D=mm$, so hat die Gleichung $tt+Duu=mm$ vier Lösungen $t,u;=m,0;\ -m,0;\ 0,1;\ 0,-1$. Sind also $F$, $f$ nur auf eine Art äquivalent, und ist eine Transformation
\[
x=\alpha x'+\beta y',\quad y=\gamma x'+\delta y'
\]
gegeben, so gibt es insgesamt \emph{vier} Transformationen,
\[
x=\pm\alpha x'\pm\beta y',\quad y=\pm\gamma x'\pm\delta y'
\]
\[
\begin{aligned}
x&=\mp\frac{\alpha B+\gamma C}{m}x'\mp\frac{\beta B+\delta C}{m}y',\\
y&=\pm\frac{\alpha A+\gamma B}{m}x'\pm\frac{\beta A+\delta B}{m}y'
\end{aligned}
\]

Sind dagegen $F$, $f$ auf zwei Arten äquivalent, oder gibt es neben der gegebenen Transformation eine weitere, ihr ungleichartige, so liefert auch diese vier Transformationen, die den ersten vier ungleichartig sind. Damit gibt es \emph{acht} Transformationen. --- Übrigens lässt sich leicht beweisen, dass in diesem Fall $F$, $f$ tatsächlich immer auf beide Arten äquivalent sind. Denn wegen $D=mm=AC-BB$ teilt $m$ auch $B$. Die Form $(\frac{A}{m},\frac{B}{m},\frac{C}{m})$ hat die Determinante $=-1$ und ist daher der Form $(1,0,1)$ oder dieser $(-1,0,-1)$ äquivalent.
\SourceLeaf{162}
Man sieht leicht, dass dieselbe Transformation, die $(\frac{A}{m},\frac{B}{m},\frac{C}{m})$ in $(\pm1,0,\pm1)$ überführt, die Form $(A,B,C)$ in $(\pm m,0,\pm m)$ überführt, die ambig ist. Da die Form $(A,B,C)$ einer ambigen Form äquivalent ist, ist sie also jeder Form, der sie äquivalent ist, sowohl eigentlich als auch uneigentlich äquivalent.

3) Ist $\frac{4D}{mm}=3$ oder $4D=3mm$, so ist $m$ gerade, und alle Lösungen der Gleichung $tt+Duu=mm$ sind die folgenden sechs:
\[
t,u;=m,0;\;-m,0;\;\tfrac12m,1;\;-\tfrac12m,-1;\;\tfrac12m,-1;\;-\tfrac12m,1
\]
Hat man also zwei ungleichartige Transformationen der Form $F$ in $f$,
\[
\begin{aligned}
x&=\alpha x'+\beta y',& y&=\gamma x'+\delta y'\\
x&=\alpha'x'+\beta'y',& y&=\gamma'x'+\delta'y'
\end{aligned}
\]
so erhält man zwölf Transformationen, nämlich sechs, die der ersten gleichartig sind,
\[
x=\pm\alpha x'\pm\beta y',\quad y=\pm\gamma x'\pm\delta y'
\]
\[
\begin{aligned}
x&=\pm(\tfrac12\alpha-\tfrac{\alpha B+\gamma C}{m})x'\pm(\tfrac12\beta-\tfrac{\beta B+\delta C}{m})y'\\
y&=\pm(\tfrac12\gamma+\tfrac{\alpha A+\gamma B}{m})x'\pm(\tfrac12\delta+\tfrac{\beta A+\delta B}{m})y'
\end{aligned}
\]
\[
\begin{aligned}
x&=\pm(\tfrac12\alpha+\tfrac{\alpha B+\gamma C}{m})x'\pm(\tfrac12\beta+\tfrac{\beta B+\delta C}{m})y'\\
y&=\pm(\tfrac12\gamma-\tfrac{\alpha A+\gamma B}{m})x'\pm(\tfrac12\delta-\tfrac{\beta A+\delta B}{m})y'
\end{aligned}
\]
und sechs, die der zweiten gleichartig sind. Diese entstehen aus den vorstehenden, indem man für $\alpha$, $\beta$, $\gamma$, $\delta$ die Werte $\alpha'$, $\beta'$, $\gamma'$, $\delta'$ setzt.

\Needspace{11\baselineskip}
Dass in diesem Fall $F,f$ immer auf beide Arten äquivalent sind, beweisen wir so: Die Form $(\frac{2A}{m},\frac{2B}{m},\frac{2C}{m})$ hat die Determinante $=-\frac{4D}{mm}=-3$ und ist deshalb (Artikel 176) entweder der Form $(\pm1,0,\pm3)$ oder dieser $(\pm2,\pm1,\pm2)$ äquivalent. Daraus erkennt man leicht, dass die Form $(A,B,C)$ entweder der Form $(\pm\tfrac12m,0,\pm\tfrac32m)$ oder dieser $(\pm m,\tfrac12m,\pm m)$\HistoricalNote{*)}{Man kann beweisen, dass die Form $(A,B,C)$ notwendig der zweiten äquivalent ist; das ist hier aber nicht nötig.} äquivalent ist. Beide sind ambig, also ist sie jeder äquivalenten Form auf beide Arten äquivalent.

4) Nimmt man $\frac{4D}{mm}=2$ an, so ergibt sich $(\frac{2B}{m})^2=4\frac{AC}{mm}-2$ und damit $\equiv2\ (\mathrm{mod.}\ 4)$.
\SourceLeaf{163}
Da aber kein Quadrat $\equiv2\ (\mathrm{mod.}\ 4)$ sein kann, ist dieser Fall unmöglich.

5) Nimmt man $\frac{4D}{mm}=1$ an, so ergibt sich $(\frac{2B}{m})^2=4\frac{AC}{mm}-1\equiv-1\ (\mathrm{mod.}\ 4)$. Da dies unmöglich ist, kann auch dieser Fall nicht eintreten.

Da außerdem $D$ weder $=0$ noch negativ ist, gibt es keine anderen als die aufgezählten Fälle.

\Needspace{8\baselineskip}
\Article{180}
\textsc{Aufgabe.} \textit{Alle Darstellungen einer gegebenen Zahl $M$ durch die Form $axx+2bxy+cyy\ldots F$ mit negativer Determinante $-D$ finden, in denen $x,y$ teilerfremde Werte annehmen.}

\textit{Lösung.} Aus Artikel 154 folgt, dass $M$ nur dann auf die verlangte Art dargestellt werden kann, wenn $-D$ ein quadratischer Rest von $M$ ist. Man bestimmt also zuerst alle verschiedenen (\textit{das heißt} inkongruenten) Werte des Ausdrucks $\surd-D\ (\mathrm{mod.}\ M)$. Diese seien $N,-N,N',-N',N'',-N''$ usw. Um die Rechnung einfacher zu machen, kann man alle $N,N'$ usw. so wählen, dass sie nicht $>\tfrac12M$ sind. Da jede Darstellung zu einem dieser Werte gehören muss, betrachten wir die Werte einzeln.

Sind die Formen $F,(M,N,\frac{D+NN}{M})$ nicht eigentlich äquivalent, so gibt es keine Darstellung von $M$, die zum Wert $N$ gehört (Artikel 168). Sind sie es dagegen, so sucht man eine eigentliche Transformation der Form $F$ in
\[
Mx'x'+2Nx'y'+\frac{D+NN}{M}y'y'
\]
und zwar sei diese
\[
x=\alpha x'+\beta y',\quad y=\gamma x'+\delta y'
\]
Dann ist $x=\alpha,y=\gamma$ eine Darstellung der Zahl $M$ durch $F$, die zu $N$ gehört. Sei der größte gemeinsame Teiler der Zahlen $A,2B,C,=m$; wir unterscheiden drei Fälle (voriger Artikel):

1) Ist $\frac{4D}{mm}>4$, so gibt es keine anderen Darstellungen, die zu $N$ gehören, als diese \textit{beiden}: $x=\alpha,y=\gamma$; $x=-\alpha,y=-\gamma$ (Artikel 169, 179).

2) Ist $\frac{4D}{mm}=4$, so gibt es \textit{vier} Darstellungen
\[
x=\pm\alpha,\quad y=\pm\gamma;\quad x=\mp\frac{\alpha B+\gamma C}{m},\quad y=\pm\frac{\alpha A+\gamma B}{m}
\]
3) Ist $\frac{4D}{mm}=3$, so gibt es \textit{sechs} Darstellungen
\SourceLeaf{164}
\[
\begin{aligned}
x&=\pm\alpha,&y&=\pm\gamma;\\
x&=\pm(\tfrac12\alpha-\tfrac{\alpha B+\gamma C}{m}),&y&=\pm(\tfrac12\gamma+\tfrac{\alpha A+\gamma B}{m})\\
x&=\pm(\tfrac12\alpha+\tfrac{\alpha B+\gamma C}{m}),&y&=\pm(\tfrac12\gamma-\tfrac{\alpha A+\gamma B}{m})
\end{aligned}
\]
Ebenso sind die Darstellungen zu bestimmen, die zu den Werten $-N.N',-N'$ usw. gehören.

\Needspace{8\baselineskip}
\Article{181}
Die Suche nach den Darstellungen der Zahl $M$ durch die Form $F$, bei denen $x,y$ keine teilerfremden Werte haben, lässt sich leicht auf den schon betrachteten Fall zurückführen. Eine solche Darstellung entstehe durch $x=\mu e,y=\mu f$, wobei $\mu$ der größte gemeinsame Teiler von $\mu e,\mu f$ ist, das heißt $e,f$ teilerfremd sind. Dann gilt $M=\mu\mu(Aee+2Bef+Cff)$, also ist die Zahl durch $\mu\mu$ teilbar. Die Substitution $x=e,y=f$ stellt dagegen die Zahl $\frac{M}{\mu\mu}$ durch die Form $F$ dar, wobei $x,y$ teilerfremde Werte haben. Ist also $M$ durch kein Quadrat außer 1 teilbar, \textit{zum Beispiel}, wenn es eine Primzahl ist, so gibt es keine solchen Darstellungen von $M$. Hat dagegen $M$ quadratische Teiler, so seien diese $\mu\mu,\nu\nu,\pi\pi$ usw. Man sucht zuerst alle Darstellungen der Zahl $\frac{M}{\mu\mu}$ durch die Form $(A,B,C)$, bei denen $x,y$ teilerfremde Werte haben. Multipliziert man diese Werte mit $\mu$, so erhält man alle Darstellungen von $M$, bei denen der größte gemeinsame Teiler der Zahlen $x,y$ gleich $\mu$ ist. Ebenso liefern alle Darstellungen von $\frac{M}{\nu\nu}$, bei denen die Werte von $x,y$ teilerfremd sind, alle Darstellungen von $M$, bei denen der größte gemeinsame Teiler der Werte von $x,y$ gleich $\nu$ ist, usw.

Damit ist klar, dass sich mit den vorigen Regeln alle Darstellungen einer gegebenen Zahl durch eine gegebene Form mit negativer Determinante finden lassen.
